% page 547 du fac-simile (image p-546 du scan)
% bloc 162.6 x 250.6 mm ; marge gauche 39.4 mm
\pgc{17.780mm}{0.000mm}{
\l{\cel{51}539}
\l{smalto.\cel{1}\sou{(tek}n.) Vitro blua, obtenat}
\l{\cel{3}vitrigebla kun kobalto. - DEFIRS.}
\l{}
\l{smeraldo. Lapido diafana, verda. -}
\l{}
\l{}
\l{smerilo. Petro harda, ek alumino, silico e fer-oxido, pulverigi}
\l{\cel{3}por polisar la lapidi, la petri, la metali, kristalo. - F}
\l{}
\l{\sou{snapa}r. (trans.) Sizar bruske, per un maxil-ago. DE}
\l{}
\l{sniflar. (trans.) Aspirar ulo per sua nazo (tabako pulverigita,}
\l{\cel{3}aquo, la humoro qua stagnas en la naztrui). - E}
\l{}
\l{snortar. (netrans.) (Dicesas precipue pri kavalo). Sniflar tr}
\l{\cel{3}bruisoze por manifestar sua repugneso ; ekpulsar aero per su}
\l{\cel{3}naztrui tre bruisoze por manifestar sua pavoro, nepacienteso}
\l{\cel{3}e c. -}
\l{}
\l{\sou{sobr}a. I. Qua su moderas pri manjo e drinko. - II.\cel{1}\sou{(met}}
\l{\cel{3}Qua esas moderata pri sua\cel{1}\sou{agi}, ne exajerema. (Ex.: Lu e}
\l{\cel{3}sobra pri komplimentar). - EFIS}
\l{}
\l{socio. Ensemblo,duroza, ek homi qui vivas obedie di legi komu}
\l{\cel{3}\sou{(an}ke aludante ula sorti de animali qui vivas ensemble : ab}
\l{\cel{3}formiki, e c.) - DEFIRS}
\l{}
\l{sociologo. Ciencisto qua su okupas pri sociologio.\cel{2}DEFIR}
\l{}
\l{sociologio. Cienco di qua la studiajo esas la cirkonstanci e}
\l{\cel{3}existo-kondicioni di la socio homara. - DEFIRS}
\l{}
\l{sociso. Peco de la intestino di porko, di mutono, quan onu p}
\l{\cel{3}igis ye karno de porko, o de bovo, tre spicizita, triturita}
\l{\cel{3}pistita, e tre presita. - EFIR}
\l{}
\l{sodo. (bot.)\cel{3}Planto ek la\cel{1}\sou{famil}io \textquotedbl{}kenopodiacei\textquotedbl{}, propra ye l}
\l{\cel{3}regioni klimato-moderata e subtropika ; un-yar-dura ; di qu}
\l{\cel{3}la folii esas quaze pulpoza, lineala, e la flori mikra e ver}
\l{\cel{3}datra ; de qui onu extraktis olim natroxo. - L. salsola sod}
\l{\cel{3}e salsola kali. - (kemio) Salo alkalioza quan onu extrakta}
\l{\cel{3}de ta planto ed anke ek la fuki.\cel{1}\sou{Simbolo kemial}a : CO3 Na}
\l{}
\l{\cel{3}DEFIR}
\l{}
\l{sodomiar.\cel{1}\sou{(netrans}.) Agar koito anusala,}
\l{\cel{3}viro e muliero, o kun best}
\l{}
\l{sofao. Sorto di repozo-lit}
\l{\cel{3}sidilo.- DEFIS}
\l{}
\l{sofismo. Argumento trompera, sedu}
\l{}
\l{sofisto. I. (en la\cel{1}\sou{epoki antiq}ua, en Grekia). Sorto de profesoro}
\l{\cel{3}di filozofio, di retoriko, qua instruktis pri quale pledar}
\l{\cel{3}e \textquotedbl{}kontre\textquotedbl{} alterne, un tezo, sen egardo pri olua exakteso.-}
\l{\cel{3}II. (nun) Persono qua argumentas trom}
}
